% Triángulo de mediación frente a distinción unilateral. Arriba, una síntesis que
% descansa en dos términos que se determinan mutuamente (`reciproca rayada`); abajo, un
% Uno que determina sin ser determinado (`unilateral`: UNA sola punta, jamás doble).
\documentclass[tikz,border=12pt]{standalone}
\usepackage{diagrama}
\begin{document}
\begin{tikzpicture}[every node/.append style={font=\px{30}}]
  \node[font=\bfseries\px{34}] at (2.5,5.6) {Mediación};
  \node (n1) at (2.5,4.6) {Síntesis};
  \node (n2) at (0.6,2.8) {Término 1};
  \node (n3) at (4.4,2.8) {Término 2};
  \draw[reciproca rayada] (n2) -- (n1);                                 % e1
  \draw[reciproca rayada] (n1) -- (n3);                                 % e2
  \draw[reciproca rayada] (n2) -- node[rotulo] {\textit{Decisión}} (n3);% e3
  \node[font=\bfseries\px{34}] at (9.5,5.6) {Distinción unilateral};
  \node (n4) at (8.2,3.4) {Uno};
  \node (n5) at (11,3.4) {Dualidad};
  \draw[unilateral] (n4) -- node[rotulo, above=1pt] {\textit{Causa unilateral}} (n5);  % e4
\end{tikzpicture}
\end{document}
